%% Career-Ops Cover Letter Template
%% Visual style matches templates/cv-template.tex (accent color, lmodern, ATS-friendly).
%% Build: pdflatex {file}.tex (or tectonic {file}.tex)

\documentclass[letterpaper,11pt]{article}

\usepackage[margin=0.85in]{geometry}
\usepackage[T1]{fontenc}
\usepackage[utf8]{inputenc}
\usepackage{lmodern}
\usepackage{microtype}
\usepackage{hyperref}
\usepackage{xcolor}
\usepackage{parskip}

\definecolor{accent}{HTML}{1F3A5F}

\hypersetup{
  colorlinks=true,
  urlcolor=accent,
  linkcolor=accent,
  pdftitle={ Ahnaf An Nafee - Cover Letter - Deepgram },
  pdfauthor={ Ahnaf An Nafee },
  pdfsubject={Cover Letter}
}

\input{glyphtounicode}
\pdfgentounicode=1

\pagestyle{empty}
\setlength{\parindent}{0pt}
\setlength{\parskip}{8pt}

\begin{document}

% --- Header (matches CV) ----------------------------------------------------
\begin{center}
  {\LARGE\bfseries\color{accent} Ahnaf An Nafee }\\[3pt]
  {\small Software Engineer \textbar{} AI Researcher \textbar{} DevOps \& Cloud Infrastructure }\\[4pt]
  \small
  \href{mailto:ahnafnafee@gmail.com}{ahnafnafee@gmail.com} \textbar{}
  540-252-8738 \textbar{}
  Fairfax, VA\\
  \href{https://linkedin.com/in/ahnafnafee}{linkedin.com/in/ahnafnafee} \textbar{}
  \href{https://github.com/ahnafnafee}{github.com/ahnafnafee} \textbar{}
  \href{https://www.ahnafnafee.dev}{ahnafnafee.dev}
\end{center}

\vspace{20pt}

% --- Date -------------------------------------------------------------------
April 27, 2026

% --- Recipient block (multi-line; use \\ between lines) ---------------------
Kris Efland\\Vice President of Engineering\\Deepgram

% --- Subject ----------------------------------------------------------------
\textbf{Re: Site Reliability Engineer \textendash{} AI \& ML Infrastructure (Kubernetes, AWS \& Terraform)}

% --- Greeting ---------------------------------------------------------------
Dear Kris,

% --- Body (3 paragraphs separated by blank lines) ---------------------------
Deepgram's hybrid AWS plus bare-metal Slurm-on-Kubernetes GPU stack is one of the more interesting reliability problems in voice AI today, and the \$130M Series C makes it clear the Nova and Aura platforms are scaling fast enough that the SRE function has to scale ahead of them. I'm writing to apply for the Site Reliability Engineer \textendash{} AI \& ML Infrastructure role because the title's three named technologies map directly onto what I've shipped in production: Kubernetes/OpenShift at enterprise scale, AWS architecture at meaningful spend, and the IaC and CI/CD discipline that keeps both honest. After 2.5 years as a DevOps Engineer at Paychex and a year as CTO at Dynasty 11 Studios, I'm now a CS PhD researcher at George Mason's DCXR Lab \textendash{} which gives me a working vocabulary for the model side of the stack you're keeping online.

At Paychex, I architected and led a cross-team standardization of Kubernetes/OpenShift deployment dictionaries across 10+ service teams, eliminating configuration conflicts that had been blocking developer velocity. I then partnered with security and infrastructure to roll out RBAC policies in OpenShift, hardening the platform against privilege-escalation risk \textendash{} the same control plane discipline a multi-tenant GPU cluster needs. The work I'm proudest of is automating the end-to-end TLS certificate lifecycle, which eliminated 100\% of certificate-expiration-related downtime; pair that with a Gradle plugin for container image certification that cut CI/CD pipeline congestion and lowered operational expense by 10\% annually, and the throughline is reliability-as-code. Earlier, as CTO at Dynasty 11, I led a monolith-to-microservices migration on AWS that cut application load and cloud spend by 80\%, and introduced GitHub Actions and Maven CI/CD that reduced manual release effort by 85\%. I haven't shipped Terraform specifically in production, but the IaC patterns from those CI/CD and AWS migration efforts transfer cleanly, and ramp time is short.

The piece that makes this role uncommon for me is that my PhD work at the DCXR Lab \textendash{} generative AI for 3D content, ML-driven graphics pipelines, and HCI in immersive XR \textendash{} means I read AI/ML infrastructure as a research-engineer, not a passive operator: I know why a bare-metal Slurm scheduler matters for training throughput, why real-time inference SLOs are different from batch, and why an SRE on a voice AI team is a force multiplier when they speak the model team's language. I'd welcome the chance to discuss how I'd ramp on the team and where I could be most useful in the first 90 days.

% --- Closing ----------------------------------------------------------------
Best regards,

\vspace{18pt}
Ahnaf An Nafee

\end{document}
